\section{Session 管理与常用命令}
%% ============================================================

\subsection{Session 的本质}

Codex 的 session 不只是聊天记录，而是随时间累积上下文、决策和操作的\textbf{工作线程}。良好的 session 管理对输出质量有重要影响。

\subsection{常用 Slash 命令}

\begin{table}[H]
\centering
\begin{tabular}{ll}
\toprule
\textbf{命令} & \textbf{功能} \\
\midrule
\texttt{/plan} & 切换到规划模式 \\
\texttt{/review} & 代码审查 \\
\texttt{/init} & 生成 starter AGENTS.md \\
\texttt{/experimental} & 切换实验性功能 \\
\texttt{/resume} & 恢复保存的会话 \\
\texttt{/fork} & 创建新线程并保留原始记录 \\
\texttt{/compact} & 压缩/总结较早的上下文 \\
\texttt{/agent} & 切换并行 agent 线程 \\
\texttt{/theme} & 选择语法高亮主题 \\
\texttt{/apps} & 在 Codex 中使用 ChatGPT apps \\
\texttt{/status} & 查看当前 session 状态 \\
\bottomrule
\end{tabular}
\caption{Codex 常用 Slash 命令一览}
\end{table}

\begin{importantbox}{一任务一线程}
保持\textbf{一个线程对应一个连贯的工作单元}。只有当工作真正产生分支时才使用 \texttt{/fork}。使用一个线程管理整个项目（而非单个任务）会导致上下文膨胀，降低输出质量。
\end{importantbox}

\begin{knowledgebox}{Subagent 工作流}
可以使用 Codex 的 subagent 功能将有限的工作从主线程卸载出去。主 agent 保持聚焦于核心问题，subagent 处理探索、测试或分类等辅助任务。
\end{knowledgebox}

\subsection{本章小结}

Session 是积累上下文的工作线程。通过 slash 命令管理线程、压缩上下文、分叉和切换 agent。坚持"一任务一线程"的原则，使用 subagent 卸载辅助工作。

%% ============================================================
